\documentclass[sigconf,nonacm]{acmart}

\settopmatter{printacmref=false}

\usepackage{amsmath}
\usepackage{amssymb}
\usepackage{booktabs}
\usepackage{graphicx}
\usepackage{url}
\usepackage{microtype}
\usepackage{array}

\graphicspath{{figures/}{results/figures/}}

\newcommand{\missingfigure}[2]{%
  \fbox{%
    \begin{minipage}[c][0.25\textheight][c]{0.92\linewidth}
      \centering
      Figure asset not found: \texttt{\detokenize{#1}}\\[0.5em]
      #2
    \end{minipage}%
  }%
}
\newcommand{\maybeincludegraphics}[2]{%
  \IfFileExists{#1}{%
    \includegraphics[width=\linewidth]{#1}%
  }{%
    \IfFileExists{results/#1}{%
      \includegraphics[width=\linewidth]{results/#1}%
    }{%
      \missingfigure{#1}{#2}%
    }%
  }%
}

\title{Conflict-Driven Creativity in AI Cross-Personality Rewriting}

\author{Kaixun Wang}
\affiliation{
  \institution{Southern University of Science and Technology}
  \city{Shenzhen}
  \country{China}
}

\begin{document}

\begin{abstract}
This report summarizes the currently available experimental results for AI cross-personality rewriting.
The report follows the structure of the provided \texttt{template.tex}, but replaces the wage-prediction task with a controlled rewriting experiment grounded in the project design documents.
The evidence is limited to files currently present under \texttt{results/} and \texttt{data/generated/}; no additional experiment run is assumed.
The current flattened analysis table contains 4296 rows: 2280 main-arm metric rows, 1920 mechanism rows, and 96 multihop rows, covering 60 source texts across academic, essay, narrative, and poetry genres.
The main arm uses four conditions: identity (T0), same-personality rewriting (T1), random-personality rewriting (T2), and directed cross-personality rewriting (T3).
All rows have non-missing automatic creativity and NLI entailment values, and all rows also contain merged LLM-judge fields.
The strongest result is that rewriting substantially improves automatic creativity over identity.
Within the pooled current analysis, the quadratic conflict term is negative for automatic creativity ($\beta_{d^2}=-0.0676$, $p=1.88\times10^{-5}$), with stronger negative curvature for value and coherence.
However, the pre-registered T3-only H1 check is not supported: its fitted vertex lies outside the empirical conflict range and the T3-only quadratic term is not significant.
Thus the current evidence supports a cautious claim: conflict is structured and directionally meaningful, but the strongest verified effect is a novelty--value trade-off rather than an unqualified T3 treatment advantage.
\end{abstract}

\keywords{AI writing, persona prompting, creativity evaluation, style transfer, LLM-as-judge, mixed-effects regression}

\maketitle

\section{Scope of This Report}

This report is a status-faithful analysis of the current result directory.
It uses the experimental logic from \texttt{docs/experiment\_workflow.md} and the measurement definitions from \texttt{docs/rubrics.md}, but the numerical claims are restricted to the present files:
\texttt{results/flat.csv}, \texttt{results/regression\_diagnostics.json}, and CSV tables in \texttt{results/tables/}.
The current \texttt{results/figures/} directory contains no PNG assets, so all figure panels are written with safe placeholders that will automatically render real plots if the expected files are regenerated.

\paragraph{Important scope note.}
The current \texttt{data/generated/main\_metrics.jsonl} contains 2280 rows and all rows have \texttt{target\_sampling=discrete}.
Therefore this report should be read as the report for the currently analysed discrete-persona result set, not as a completed continuous-target report.

\section{Experimental Design}

\subsection{Persona Space and Conflict}

The project defines an interpretable two-dimensional Space-H.
The first axis, $S$, represents a rational--emotional contrast inspired by dual-process theory.
The second axis, $R$, represents a conservative--adventurous register-risk contrast.
Each source text receives a coordinate in this plane, and each target persona is also represented by a coordinate.
Conflict intensity is the scalar distance $d_H$ between source and target, while directional conflict is $\Delta_H=(\Delta S,\Delta R)$.

\subsection{Four-Arm Main Experiment}

The main arm follows the quasi-causal structure described in the workflow document.
T0 is the identity baseline.
T1 rewrites the text using the nearest persona and estimates the pure rewrite effect.
T2 rewrites under a random persona and acts as a placebo for arbitrary stylistic change.
T3 rewrites toward a directed cross-personality target.
The current main metric file has 60 sources, two generator roles, and 2280 analysed rows.

\begin{table}[t]
\centering
\small
\begin{tabular}{lrrrr}
\toprule
Source file & Rows & T0 & T1/T2 & T3 or modes \\
\midrule
\texttt{main\_metrics} & 2280 & 120 & 360 / 360 & 1440 T3 \\
\texttt{mechanism\_metrics} & 1920 & -- & -- & 480 each M0--M3 \\
\texttt{multihop\_metrics} & 96 & -- & -- & 96 hop rows \\
\midrule
Total & 4296 & 120 & 720 & 3456 non-identity rows \\
\bottomrule
\end{tabular}
\caption{Current flattened result coverage from \texttt{results/flat.csv}. Multihop rows do not use T0--T3 condition labels in the flattened table.}
\label{tab:coverage}
\end{table}

\section{Measurement Model}

The primary automatic dependent variable is
\[
  C_{\mathrm{auto}} = N_{\mathrm{auto}} \times V_{\mathrm{auto}}.
\]
The novelty term combines genre-normalised distinct-2, novel n-gram ratio, and embedding-distance signals when available.
The value term is the arithmetic mean of an NLI entailment proxy and a bounded coherence score derived from perplexity:
\[
  V_{\mathrm{auto}} = \frac{Fidelity_{\mathrm{NLI}} + Coherence_{\mathrm{auto}}}{2}.
\]
The project documentation explicitly treats LLM judge scores as validity and sensitivity evidence, not as a replacement for the automatic primary outcome.
In the current run, judge fields are merged for all 4296 rows.

\begin{table}[t]
\centering
\small
\begin{tabular}{lrrrrr}
\toprule
Condition & Creativity & Novelty & Value & NLI & Coherence \\
\midrule
T0 & 0.2317 & 0.3334 & 0.7044 & 0.7486 & 0.6602 \\
T1 & 0.3191 & 0.4041 & 0.7918 & 0.8066 & 0.7770 \\
T2 & 0.3097 & 0.4186 & 0.7418 & 0.7143 & 0.7693 \\
T3 & 0.3124 & 0.4232 & 0.7409 & 0.7144 & 0.7675 \\
\midrule
M0 & 0.2937 & 0.4136 & 0.7143 & 0.6230 & 0.8055 \\
M1 & 0.2937 & 0.4167 & 0.7086 & 0.6100 & 0.8073 \\
M2 & 0.3048 & 0.4371 & 0.7019 & 0.5993 & 0.8046 \\
M3 & 0.2926 & 0.4049 & 0.7272 & 0.6498 & 0.8045 \\
\bottomrule
\end{tabular}
\caption{Mean automatic metrics by condition from \texttt{results/flat.csv}. T1 has the highest main-arm mean creativity; T3 is slightly above T2 but below T1 in aggregate.}
\label{tab:condition}
\end{table}

\section{Descriptive Results}

\subsection{Genre-Level Patterns}

Genre differences are substantial even after genre-normalised novelty is used.
Essay and poetry have the highest mean creativity in the current flattened table, while narrative is lowest.
This is not solely a novelty effect: value, NLI, and coherence differ sharply by genre as well.

\begin{table}[t]
\centering
\small
\begin{tabular}{lrrrrr}
\toprule
Genre & Creativity & Novelty & Value & NLI & Coherence \\
\midrule
Academic & 0.2778 & 0.4119 & 0.6832 & 0.5958 & 0.7705 \\
Essay & 0.3503 & 0.4572 & 0.7702 & 0.7657 & 0.7747 \\
Narrative & 0.2635 & 0.3698 & 0.7151 & 0.5318 & 0.8985 \\
Poetry & 0.3405 & 0.4377 & 0.7838 & 0.8996 & 0.6679 \\
\bottomrule
\end{tabular}
\caption{Genre-level metric means from \texttt{results/flat.csv}. Poetry has high NLI but lower coherence; narrative has the lowest creativity and NLI in the current run.}
\label{tab:genre}
\end{table}

\begin{figure*}[t]
\centering
\begin{minipage}[t]{0.49\textwidth}
\centering
\maybeincludegraphics{figures/inverted_u_creativity_auto.png}{Regenerate by running \texttt{python -m src.analyze --metric creativity\_auto}.}
\end{minipage}\hfill
\begin{minipage}[t]{0.49\textwidth}
\centering
\maybeincludegraphics{figures/directional_field_creativity_auto.png}{Regenerate by running \texttt{python -m src.analyze --metric creativity\_auto}.}
\end{minipage}
\caption{Expected analysis figures. The current result directory contains no PNG files, so placeholders are shown until figures are regenerated.}
\label{fig:expected}
\end{figure*}

\section{Conflict Intensity}

\subsection{Pooled Quadratic Results}

The pooled current regression over all loaded metric files shows a negative quadratic term for automatic creativity.
The same sign is robust for geometric creativity and for the winsorised creativity sensitivity run.
The decomposition indicates that this curvature is mainly associated with value and coherence loss at higher conflict; novelty alone does not show a significant negative quadratic in the non-winsorised model.

\begin{table}[t]
\centering
\small
\begin{tabular}{lrrr}
\toprule
Metric & $n$ & $\beta_{d^2}$ & $p(d^2)$ \\
\midrule
Creativity auto & 4200 & -0.0676 & $1.88{\times}10^{-5}$ \\
Creativity geom & 4200 & -0.0680 & $5.77{\times}10^{-4}$ \\
Value auto & 4200 & -0.1816 & $2.10{\times}10^{-7}$ \\
Coherence auto & 4200 & -0.1608 & $3.25{\times}10^{-6}$ \\
Novelty auto & 4200 & 0.0198 & 0.1615 \\
\bottomrule
\end{tabular}
\caption{Key quadratic terms from \texttt{results/regression\_diagnostics.json}. These regressions are pooled over the current loaded metrics files, not only the pre-registered T3 main subset.}
\label{tab:quadratic}
\end{table}

\subsection{Pre-Registered H1 Check}

The pre-registered T3-only H1 diagnostic is weaker than the pooled result.
For T3 rows in the main metric file, the fitted vertex is $d=-0.516$, outside the empirical T3 conflict range.
The T3-only quadratic term is not significant ($p=0.560$), and both pre-registration falsification flags are triggered.
This means the strongest honest conclusion is not that the pre-registered T3-only inverted-U has been confirmed.
Instead, the current evidence shows a broader pooled concavity in creativity and strong value/coherence collapse at high conflict.

\begin{table}[t]
\centering
\small
\begin{tabular}{lr}
\toprule
H1 diagnostic & Value \\
\midrule
T3 main rows & 1440 \\
$\beta_d$ & -0.0167 \\
$\beta_{d^2}$ & -0.0161 \\
$p(d^2)$ & 0.5599 \\
Fitted vertex $d$ & -0.5164 \\
Empirical $d$ q03 / q97 & 0.1275 / 0.9055 \\
Vertex inside pre-reg band & False \\
Vertex inside sample mass & False \\
\bottomrule
\end{tabular}
\caption{Pre-registered H1 diagnostic from \texttt{results/tables/h1\_preregistration\_check.csv}.}
\label{tab:h1}
\end{table}

\section{Causal Contrasts}

The bucketed causal contrast table compares T3 to T1 and T2 within conflict bins.
The current evidence is small and mixed.
In the low bin, T3 is slightly below T1 but above T2.
In the mid bin, T3 is above T2 by about 0.0021.
In the high bin, T3 is below T2 by about 0.0355.
This pattern does not support a simple statement that directed conflict always improves creativity.

\begin{table}[t]
\centering
\small
\begin{tabular}{lrrrrr}
\toprule
Bucket & $n_{T1}$ & $n_{T2}$ & $n_{T3}$ & T3--T1 & T3--T2 \\
\midrule
Low & 360 & 212 & 786 & -0.0009 & 0.0075 \\
Mid & 0 & 128 & 564 & -- & 0.0021 \\
High & 0 & 20 & 90 & -- & -0.0355 \\
\bottomrule
\end{tabular}
\caption{Bucketed causal contrasts from \texttt{results/tables/causal\_creativity\_auto.csv}. Missing T1 entries reflect absent T1 rows in the corresponding conflict bucket.}
\label{tab:causal}
\end{table}

\section{Conflict Direction}

Directional regression checks whether the vector components of conflict matter beyond scalar distance.
For automatic creativity, $\Delta S$ is positive and significant, $\Delta R$ is not significant on its own, $\Delta R^2$ is significantly negative, and the interaction $\Delta S\Delta R$ is strongly positive.
The absolute magnitude term is not significant after the directional terms enter.
This suggests that the geometry of conflict is more informative than distance alone.

\begin{table}[t]
\centering
\small
\begin{tabular}{lrr}
\toprule
Term & Coefficient & $p$ \\
\midrule
$\Delta S$ & 0.0065 & $3.15{\times}10^{-8}$ \\
$\Delta R$ & 0.0008 & 0.3951 \\
$\Delta S^2$ & -0.0006 & 0.8067 \\
$\Delta R^2$ & -0.0051 & 0.0194 \\
$\Delta S\Delta R$ & 0.0067 & $7.50{\times}10^{-25}$ \\
$|\Delta|$ & -0.0002 & 0.9767 \\
\bottomrule
\end{tabular}
\caption{Directional creativity regression from \texttt{results/regression\_diagnostics.json}.}
\label{tab:direction}
\end{table}

\section{Mechanism and Multihop Branches}

\subsection{Mechanism Experiment}

The mechanism branch is present in the current metrics table.
At the high-conflict threshold $d=0.7906$, M2 has the highest mean automatic creativity among M0--M3.
This is consistent with the planned idea that a two-stage paraphrase-then-style pipeline can mitigate some high-conflict damage.
The result should still be treated as suggestive, because this table is a mean contrast rather than a full causal model.

\begin{table}[t]
\centering
\small
\begin{tabular}{lrrr}
\toprule
Mode & Mean creativity & Std. & Count \\
\midrule
M0 & 0.2908 & 0.0759 & 342 \\
M1 & 0.2919 & 0.0806 & 342 \\
M2 & 0.3067 & 0.0771 & 342 \\
M3 & 0.2897 & 0.0774 & 342 \\
Joint & 0.2930 & 0.0811 & 418 \\
\bottomrule
\end{tabular}
\caption{High-conflict mechanism summary from \texttt{results/tables/mechanism\_creativity\_auto.csv}.}
\label{tab:mechanism}
\end{table}

\subsection{Multihop Prediction Check}

The multihop branch has a prediction-check table, but the evidence is heterogeneous.
Essay and poetry mid-bin rows support the bridge-over-direct prediction, while academic rows do not.
The table is therefore best reported as exploratory path-dependence evidence rather than a confirmed general effect.

\begin{table}[t]
\centering
\small
\begin{tabular}{llrrr}
\toprule
Genre & Bin & Multihop & Direct T3 & Delta \\
\midrule
Academic & Low & 0.2700 & 0.2994 & -0.0295 \\
Academic & Mid & 0.2961 & 0.3004 & -0.0043 \\
Academic & High & 0.2553 & 0.2737 & -0.0183 \\
Essay & Mid & 0.4506 & 0.3517 & 0.0988 \\
Narrative & High & 0.1710 & 0.2563 & -0.0853 \\
Poetry & Mid & 0.4627 & 0.3606 & 0.1021 \\
\bottomrule
\end{tabular}
\caption{Selected multihop contrasts from \texttt{results/tables/multihop\_prediction\_check.csv}. The full CSV contains all available genre-bin rows.}
\label{tab:multihop}
\end{table}

\section{Judge Validity Evidence}

The judge--automatic correlation matrix is available in \texttt{results/tables/judge\_auto\_validity\_correlation.csv}.
Judge novelty correlates moderately with automatic novelty ($r=0.435$), but judge fidelity is negatively correlated with judge novelty ($r=-0.730$), reflecting the expected novelty--fidelity tension.
Automatic creativity correlates with automatic novelty ($r=0.611$) and value ($r=0.685$), while its direct correlation with judge novelty is near zero ($r=0.016$).
This supports the documentation's caution that LLM judge scores should be interpreted as validity evidence, not as a drop-in replacement for the primary automatic metric.

\begin{table}[t]
\centering
\small
\begin{tabular}{lrr}
\toprule
Pair & Pearson $r$ & Interpretation \\
\midrule
Judge novelty vs. auto novelty & 0.435 & moderate convergence \\
Judge fidelity vs. judge novelty & -0.730 & strong trade-off \\
Auto creativity vs. auto novelty & 0.611 & composite dependence \\
Auto creativity vs. value auto & 0.685 & composite dependence \\
Auto creativity vs. judge novelty & 0.016 & little direct alignment \\
\bottomrule
\end{tabular}
\caption{Selected validity correlations from \texttt{results/tables/judge\_auto\_validity\_correlation.csv}.}
\label{tab:judge}
\end{table}

\section{Coordinate Reliability}

The coordinate reliability file currently contains 60 sources and compares two coordinate scorers.
The summary rows report near-zero ICC values for both $S$ and $R$.
This is a serious limitation for claims that depend on the absolute reliability of source coordinates.
The analysis can still be useful as an internal geometry-based experiment, but any paper-style claim about Space-H should explicitly disclose the coordinate reliability weakness and avoid over-claiming the precision of $d_H$.

\section{Limitations}

First, the current report is based on existing result files only.
It does not rerun missing figure generation, continuous-target generation, or additional metrics.
Second, all currently analysed main rows in \texttt{main\_metrics.jsonl} use discrete target sampling; the continuous-target branch is not represented in the final metric table used here.
Third, \texttt{results/figures/} is empty, so figures in this report are reproducible placeholders rather than verified image assets.
Fourth, the H1 pre-registered T3-only check is not supported, even though pooled regressions show strong negative quadratic curvature.
Finally, coordinate reliability is weak in the current file, which should temper claims about exact Space-H distances.

\section{Reproducibility}

The report was written against the following current artefacts:
\begin{itemize}\itemsep2pt
\item \texttt{results/flat.csv}: flattened analysis table with 4296 rows.
\item \texttt{data/generated/main\_metrics.jsonl}: 2280 main-arm metric rows.
\item \texttt{data/generated/mechanism\_metrics.jsonl}: 1920 mechanism metric rows.
\item \texttt{data/generated/multihop\_metrics.jsonl}: 96 multihop metric rows.
\item \texttt{results/regression\_diagnostics.json}: pooled and directional regression diagnostics.
\item \texttt{results/tables/*.csv}: H1, causal, mechanism, multihop, judge-validity, and coordinate tables.
\end{itemize}

To regenerate missing plots from the current metrics files, run:
\begin{verbatim}
python -m src.analyze --metric creativity_auto
\end{verbatim}
If the intended final study is the continuous-target version, the required next step is to complete continuous generation, rerun metrics on that file, merge judge scores, and rerun analysis before replacing the numerical claims in this report.

\section{Conclusion}

The current results support a careful interpretation.
Rewriting clearly improves automatic creativity over identity, and conflict geometry matters.
The pooled regressions show a strong negative quadratic relationship between conflict and creativity, driven mainly by value and coherence decline at higher conflict.
However, the pre-registered T3-only inverted-U diagnostic is not supported, and T3 does not dominate T1/T2 in simple aggregate means.
The most defensible conclusion is therefore conditional: persona conflict can be productive in some regions and directions, but excessive or poorly aligned conflict damages value, and current evidence should not be framed as a universal T3 advantage.

\begin{thebibliography}{9}
\bibitem{kahneman}
D.~Kahneman.
\newblock \textit{Thinking, Fast and Slow}.
\newblock Farrar, Straus and Giroux, 2011.

\bibitem{boden}
M.~Boden.
\newblock \textit{The Creative Mind: Myths and Mechanisms}.
\newblock Routledge, 2004.

\bibitem{runco}
M.~Runco and G.~Jaeger.
\newblock The standard definition of creativity.
\newblock \textit{Creativity Research Journal} 24, 1 (2012), 92--96.

\bibitem{yerkes}
R.~Yerkes and J.~Dodson.
\newblock The relation of strength of stimulus to rapidity of habit-formation.
\newblock \textit{Journal of Comparative Neurology and Psychology} 18 (1908), 459--482.

\bibitem{shannon}
C.~Shannon.
\newblock A mathematical theory of communication.
\newblock \textit{Bell System Technical Journal} 27 (1948), 379--423.
\end{thebibliography}

\end{document}
